\begin{table}[pos=htbp]
\caption{Source preparation, learning and deployment procedure.}
\label{tab:algorithm}
\begin{minipage}{.98\linewidth}\small
\begin{enumerate}
\item Partition the observed source recordings into training and validation
roles. Select same-class training-context pairs using the declared source
metadata; retain bearing identity for PU.
\item Form the response, shaft-harmonic and stationary dictionaries. Run the
fixed numerical inverse in Eq.~\eqref{eq:inverse} for each registered source pair.
\item For each source context, render $u_1$ and the two fixed-rate $u_q$ signals.
Construct $\mathcal T_q(x)$ with Eq.~\eqref{eq:edit}, then crop the three
observation-supported queries and cache their log-FFT features.
\item For each training update, sample a balanced factual batch and same-class,
same-source-condition bank slots. Draw the shared edit mask and rate direction;
form the auxiliary features using Eq.~\eqref{eq:aux}.
\item Draw the Mixup coefficient and permutation. Update the shared classifier
using Eq.~\eqref{eq:loss} with the common $6B$ forward layout.
\item Freeze the classifier after 2000 updates for the primary evaluation.
Evaluate source-validation selection separately as a secondary policy.
\item At deployment, apply the log-FFT frontend and trained classifier to each
vibration window. Average probabilities only when reporting recording-level
diagnosis.
\end{enumerate}
\end{minipage}
\end{table}
